\documentclass[11pt]{amsart}
\usepackage[margin=1.1in]{geometry}
\usepackage{amsmath,amssymb,amsthm,mathtools,booktabs,enumitem}
\usepackage[hidelinks]{hyperref}

\newtheorem{theorem}{Theorem}[section]
\newtheorem{proposition}[theorem]{Proposition}
\newtheorem{lemma}[theorem]{Lemma}
\newtheorem{corollary}[theorem]{Corollary}
\newtheorem{conjecture}[theorem]{Conjecture}
\theoremstyle{remark}
\newtheorem{remark}[theorem]{Remark}
\theoremstyle{definition}
\newtheorem{definition}[theorem]{Definition}

\newcommand{\Z}{\mathbb Z}
\newcommand{\Q}{\mathbb Q}
\newcommand{\R}{\mathbb R}
\newcommand{\C}{\mathbb C}
\newcommand{\HH}{\mathbb H}
\newcommand{\cG}{\mathcal G}
\newcommand{\cP}{\mathcal P}
\newcommand{\cE}{\mathcal E}
\newcommand{\cO}{\mathcal O}
\newcommand{\cR}{\mathcal R}
\newcommand{\sgn}{\operatorname{sgn}}
\newcommand{\Rea}{\operatorname{Re}}
\newcommand{\Ima}{\operatorname{Im}}
\newcommand{\Log}{\operatorname{Log}}
\newcommand{\Arg}{\operatorname{Arg}}
\newcommand{\e}{\mathrm e}
\newcommand{\KNG}{K_{\mathrm{NG}}}

\title[Real powers of the G\"ollnitz--Gordon product]{Sign patterns of real powers of the G\"ollnitz--Gordon product: a complete resolution of a conjecture of Schlosser and Zhou}
\author{Jaideep Sai Padhi}
\address{Department of Computer Science, Purdue University, West Lafayette, IN, USA}
\email{jpadhi@purdue.edu}
\date{\today}

\begin{document}
\begin{abstract}
Let $Q_8(q)=(q,q^7;q^8)_\infty/(q^3,q^5;q^8)_\infty$ and write $Q_8(q)^\delta=\sum_{n\ge0}c_\delta(n)q^n$ for real $\delta$. Schlosser and Zhou conjectured that for $\beta\le\delta\le4$, where $\beta=2.66447911\ldots$ is the point at which $c_\delta(14)$ vanishes, the signs of $c_\delta(n)$ follow a fixed pattern of period $8$. We show that this holds for every $\delta\in[8/3,4]$ and all $n\ge0$, and that it fails for every $\delta\in[\beta,8/3)$. We also prove the remaining open part of the conjecture, $-1<\delta<-0.99$, where the coefficients degenerate as $\delta\to-1$, and give certified proofs of the parts previously treated by He and Li, so that the conjecture is completely resolved. The failure first occurs near $n\approx7.6\times10^5$, far beyond the range of the numerical evidence. The mechanism is that at the residues $n\equiv2,6\pmod8$ the leading circle-method term vanishes identically, so the sign is decided by two secondary terms whose exponential growth rates coincide exactly at $\delta=8/3$. The proof is self-contained. It combines an exact description of the cusp behaviour of $Q_8$, obtained from the Weil representation attached to a lattice of level $64$, with an explicit Hardy--Ramanujan--Rademacher expansion and certified computations in exact and ball arithmetic.
\end{abstract}

\maketitle

\section{Introduction}

For real $\delta$ let
\[
Q_8(q)=\frac{(q,q^7;q^8)_\infty}{(q^3,q^5;q^8)_\infty},\qquad Q_8(q)^\delta=\exp(\delta\log Q_8(q))=\sum_{n\ge0}c_\delta(n)q^n,
\]
where $\log Q_8(0)=0$. The function $q^{1/2}Q_8(q)$ is the G\"ollnitz--Gordon continued fraction. Taking logarithmic derivatives gives the recurrence
\begin{equation}\label{eq:rec}
n\,c_\delta(n)=\delta\sum_{k=1}^n b_k\,c_\delta(n-k),\qquad b_k=\sum_{d\mid k}\epsilon(d)\,d,
\end{equation}
where $\epsilon(d)=1$ for $d\equiv3,5$, $\epsilon(d)=-1$ for $d\equiv1,7\pmod 8$, and $\epsilon(d)=0$ otherwise. In particular each $c_\delta(n)$ is a polynomial in $\delta$ of degree at most $n$ with rational coefficients.

Schlosser and Zhou \cite{SZ} studied sign patterns of real powers of infinite products and posed a number of conjectures in an appendix. For the G\"ollnitz--Gordon product they conjectured the following.

\begin{conjecture}[{\cite[Conj.~21]{SZ}}]\label{conj}
The $q$-series coefficients of $Q_8(q)^\delta$ exhibit
\begin{enumerate}[label=(\alph*),leftmargin=2em]
\item for $\delta=2$ the length $16$ sign pattern $+-++-+--+--+-++-$;
\item for $2.664479110226972\ldots\approx\beta\le\delta\le4$ the sign pattern
\begin{equation}\label{eq:sigma}
\sigma=(+,-,+,+,-,+,-,-)\qquad(n\equiv0,1,\dots,7\bmod 8),
\end{equation}
where $\beta$ is the unique real root in $(2,3)$ of $x^{12}-90x^{11}+1457x^{10}+30486x^9-537081x^8+1892346x^7-3683653x^6-837509646x^5+774767020x^4+3333687384x^3-40887173664x^2+94379731200x+49816166400$;
\item for $-1<\delta\le\frac{7-\sqrt{73}}2\approx-0.77200187$ the sign pattern $\sigma'=(+,+,+,-,-,-,-,+)$;
\item for $\delta=-2$ the length $16$ sign pattern $+++++----+++----$.
\end{enumerate}
\end{conjecture}

As in \cite{SZ}, where the analogous Conjecture~2 is phrased as non-negativity of the coefficients of a dissection, we read a sign pattern as allowing zero coefficients. Exactly three zeros occur in the ranges above: $c_4(3)=0$, $c_{-2}(8)=0$, and $c_{b}(4)=0$ at the endpoint $b=\frac{7-\sqrt{73}}2$ (which is where the range in (c) comes from).

Parts (a), (d), and part (c) for $-0.99\le\delta\le\frac{7-\sqrt{73}}2$ were proved by He and Li \cite{HL}, with finite checks done in floating point. We resolve what remained, and give certified proofs of all parts.

\begin{theorem}\label{thm:A}
For every $\delta\in[8/3,4]$ and every $n\ge0$, $c_\delta(n)$ has the sign $\sigma(n\bmod 8)$. The only vanishing coefficient is $c_4(3)=0$.
\end{theorem}

\begin{theorem}\label{thm:B}
Let $2<\delta<8/3$. There is $n_1(\delta)$ such that $c_\delta(n)<0$ for all $n\ge n_1(\delta)$ with $n\equiv10\pmod{16}$, and $c_\delta(n)>0$ for all $n\ge n_1(\delta)$ with $n\equiv14\pmod{16}$. For example,
\[
c_{2.66448}(1\,000\,010)<-5.7\cdot10^{387}.
\]
\end{theorem}

\begin{theorem}\label{thm:C}
Parts (a), (c) and (d) of Conjecture~\ref{conj} hold. In particular, for every $\delta\in(-1,\frac{7-\sqrt{73}}2]$ and every $n\ge0$, $c_\delta(n)$ has the sign $\sigma'(n\bmod8)$.
\end{theorem}

\begin{corollary}
Conjecture~\ref{conj} is completely resolved: parts (a), (c), (d) are true; part (b) is true for $\delta\in[8/3,4]$ and false for every $\delta\in[\beta,8/3)$.
\end{corollary}

For $\delta=\beta$ the first failures predicted by the asymptotic expansion are at $n=759\,390$ (residue $14\bmod16$) and $n=759\,402$ (residue $10\bmod16$). The crossover point increases rapidly as $\delta\to8/3^-$; for instance it is about $8.3\times10^{10}$ at $\delta=2.66666$.

\subsection*{The mechanism}
The dominant cusps of $Q_8$ are $3/8$ and $5/8$. Their contributions combine to $2\cos(3\pi n/4)$ times a Bessel function, which vanishes \emph{identically} for $n\equiv2,6\pmod 8$. For those residues the sign of $c_\delta(n)$ is decided by two lower-order terms:
\begin{itemize}[leftmargin=2em]
\item a secondary term at the cusps $3/8,5/8$, of size $\exp\big(\tfrac\pi4\sqrt{(\delta-2)(2n+\delta)}\big)$, which always has the sign required by \eqref{eq:sigma};
\item the leading term at the cusps of denominator $16$, of size $\exp\big(\tfrac\pi8\sqrt{\delta(2n+\delta)}\big)$, whose sign in the classes $n\equiv10,14\pmod{16}$ is wrong when $2<\delta<3$.
\end{itemize}
The two rates agree exactly when $\delta-2=\delta/4$, that is $\delta=8/3$. At $\delta=8/3$ the correct term still wins by a fixed factor $1/0.375$. Because the two rates differ by only about $10^{-3}\sqrt n$ near $\beta$, the failure is invisible for small $n$.

\subsection*{Method}
The proof is self-contained apart from classical facts: Poisson summation, the generation of $SL_2(\Z)$ by $S$ and $T$, Ford circles, and integral representations of $I_1$.
\begin{itemize}[leftmargin=2em]
\item \S\ref{sec:prelim} writes $Q_8$ as a quotient of theta series.
\item \S\ref{sec:cusps} determines the behaviour of $Q_8$ at every cusp. The pair of theta series lies in a finite projective orbit of $48$ elements under the Weil representation of $SL_2(\Z)$ at level $64$; exactly one of its ten cusp types grows, and it has a rigid shape (Lemma~\ref{lem:growing}).
\item \S\ref{sec:circle} carries out Rademacher's method with fully explicit error terms.
\item \S\ref{sec:A} proves Theorem~\ref{thm:A}, splitting into $n\le700$ (exact polynomial arithmetic), $701\le n\le20000$ (ball arithmetic), and $n>20000$ (an analytic tail estimate).
\item \S\ref{sec:B} proves Theorem~\ref{thm:B}.
\item \S\ref{sec:computations} describes the computations and what they rely on.
\end{itemize}

\subsection*{Relation to earlier work}
For $\delta=1$ the sign pattern of the coefficients of the G\"ollnitz--Gordon continued fraction was established by Richmond and Szekeres \cite{RS} using Iseki's transformation formulae \cite{Iseki}, and Hirschhorn \cite{Hirschhorn} studied the expansion further. Asymptotics for coefficients of eta-quotients and related products were developed by Chern \cite{Chern}. Several of the conjectures in \cite{SZ} have since been proved; see \cite{HL} and \cite{BHK}.
He and Li \cite{HL} developed a general explicit circle method for real powers of such products and proved several of the conjectures in \cite{SZ}. For Conjecture~\ref{conj} their method applies at $\delta=2$ and for some negative $\delta$, but a technical condition fails for $\beta\le\delta\le4$. We do not use results of \cite{HL}; the cusp analysis of \S\ref{sec:cusps} replaces the general transformation formula used there.

\section{Theta-quotient form}\label{sec:prelim}

Throughout, $e(x)=\e^{2\pi ix}$ and $\zeta_m=e(1/m)$. By the Jacobi triple product with base $q^8$,
\[
(q^a,q^{8-a},q^8;q^8)_\infty=\sum_{m\in\Z}(-1)^mq^{4m^2+(a-4)m}.
\]
The factor $(q^8;q^8)_\infty$ cancels in $Q_8$. Putting $j=8m+a-4$ and $q=e(\tau)$ we obtain
\begin{equation}\label{eq:theta-quotient}
Q_8=e(-\tau/2)\,\frac{S_1(\tau)}{S_3(\tau)},\qquad S_a(\tau)=\sum_{j\in\Z}\psi_a(j)\,e(\tau j^2/16),
\end{equation}
where $\psi_a$ is $16$-periodic, with $\psi_a(j)=(-1)^{(j-a+4)/8}$ if $j\equiv a-4\pmod8$ and $\psi_a(j)=0$ otherwise.

\medskip\noindent\textbf{Notation.} For coprime $h,k$ with $k\ge1$ we write
\[
\tau=\frac hk+\frac{iz}{k^2}\qquad(\Rea z>0),\qquad M=16k .
\]

\begin{lemma}\label{lem:poisson}
With the principal square root,
\[
S_a\Big(\frac hk+\frac{iz}{k^2}\Big)=\frac1{16}\sqrt{\frac8z}\,\sum_{\ell\in\Z}\cG_a(\ell)\,X^{\ell^2},\qquad X=\e^{-\pi/(32z)},\qquad \cG_a(\ell)=\sum_{s\bmod M}\psi_a(s)\,\zeta_M^{hs^2+\ell s}.
\]
\end{lemma}

\begin{proof}
We have $e(\tau j^2/16)=\zeta_M^{hj^2}\e^{-\pi zj^2/(8k^2)}$, and $j\mapsto\psi_a(j)\zeta_M^{hj^2}$ is $M$-periodic. Writing $j=s+Mt$ and applying Poisson summation in $t$ to the Gaussian gives
\[
\sum_{t}\e^{-\pi z(s+Mt)^2/(8k^2)}=\frac1M\sqrt{\frac{8k^2}z}\sum_\ell e(\ell s/M)\,\e^{-\pi\ell^2/(32z)}.
\]
All series converge absolutely and are holomorphic in $z$.
\end{proof}

\section{Cusp structure}\label{sec:cusps}

For coprime $h,k$ choose $\gamma=\begin{psmallmatrix}h&b\\k&d\end{psmallmatrix}\in SL_2(\Z)$. With $\tau=h/k+iz/k^2$ a direct computation gives
\begin{equation}\label{eq:tauprime}
\tau'=\gamma^{-1}\tau=-\frac dk+\frac iz,\qquad \Ima\tau'=\Rea(1/z).
\end{equation}

\subsection{Theta functions of level 64}
For $c\in\Z/32\Z$ let $\vartheta_c(\tau)=\sum_{J\equiv c\,(32)}e(\tau J^2/64)$. Then $\vartheta_c=\vartheta_{-c}$, and comparing with \eqref{eq:theta-quotient},
\begin{equation}\label{eq:S-theta}
S_1=\vartheta_6-\vartheta_{10},\qquad S_3=\vartheta_2-\vartheta_{14}.
\end{equation}

\begin{lemma}\label{lem:weil}
We have $\vartheta_c(\tau+1)=e(c^2/64)\,\vartheta_c(\tau)$ and
\[
\vartheta_c(-1/\tau)=\frac{(i/(32\tau))^{-1/2}}{32}\sum_{d\bmod 32}e(cd/32)\,\vartheta_d(\tau).
\]
\end{lemma}

\begin{proof}
The first identity is clear. For the second, apply Poisson summation to $m\mapsto\e^{\pi i\tau(32m+c)^2/32}$. Its Fourier transform at $\nu$ is $\frac1{32}(-i\tau/32)^{-1/2}e(\nu c/32)\,\e^{-\pi i\nu^2/(32\tau)}$. Grouping $\nu$ by residues modulo $32$ and replacing $\tau$ by $-1/\tau$ gives the claim.
\end{proof}

For a symmetric vector $v=(v_c)_{c\in\Z/32}$ (so $v_c=v_{-c}$) put $v\cdot\vartheta=\sum_cv_c\vartheta_c$. Let
\[
D=\operatorname{diag}\big(\zeta_{128}^{2c^2}\big),\qquad E=\big(\zeta_{128}^{4cd}\big)_{c,d}.
\]
By Lemma~\ref{lem:weil},
\[
(v\cdot\vartheta)(T\tau)=(vD)\cdot\vartheta(\tau),\qquad (v\cdot\vartheta)(S\tau)=s(\tau)\,(vE)\cdot\vartheta(\tau),
\]
where the scalar $s(\tau)$ does not depend on $v$. Since $SL_2(\Z)$ is generated by $S$ and $T$, and $S^{-1}=S^3$ while $T^{-1}$ acts projectively as $D^{127}$ (because $D^{128}=I$), every $\gamma$ is a word $g_1\cdots g_r$ in $S$ and $T$ without inverses, up to the central element $-I$, which acts on symmetric vectors by a scalar. Iterating,
\begin{equation}\label{eq:R}
\frac{S_1}{S_3}(\gamma\tau')=R_{(w_1,w_3)}(\tau'):=\frac{w_1\cdot\vartheta(\tau')}{w_3\cdot\vartheta(\tau')},\qquad (w_1,w_3)=(v_1,v_3)\rho(g_1)\cdots\rho(g_r),
\end{equation}
where $v_1,v_3$ are the coefficient vectors in \eqref{eq:S-theta} and $\rho(S)=E$, $\rho(T)=D$. The accumulated scalar is the same for both components, so $R$ depends only on the projective class of the pair.

\begin{proposition}\label{prop:orbit}
The projective orbit $\cO$ of $(v_1,v_3)$ under right multiplication by $D$ and $E$ has exactly $48$ elements and is closed under both maps. Under $D$ it decomposes into $10$ cycles, which we call \emph{cusp types}.
\end{proposition}

\begin{proof}
This is an exact computation in $\Z[\zeta_{128}]=\Z[x]/(x^{64}+1)$. Starting from $(v_1,v_3)$, images under $D$ and $E$ were generated. Each image was tested for projective equality with the known elements by checking all cross-products $p_iq_j=p_jq_i$ exactly. All $96$ images coincide with one of $48$ elements.
\end{proof}

Near $\Ima\tau'=\infty$,
\[
w\cdot\vartheta(\tau')=\sum_{J\ge0}A_J\,q'^{J^2/64},\qquad q'=e(\tau'),\quad A_0=w_0,\quad A_J=2w_{J\bmod32}\ (J\ge1).
\]
So the leading exponent of the $a$-th component is $J_a=\min\{J\ge0:w_{a,J}\ne0\}$. Since $w$ is symmetric, no cancellation between $J$ and $-J$ can occur.

\begin{proposition}\label{prop:types}
Exactly one cusp type is \emph{growing}, i.e.\ has $J_1<J_3$, and for it $(J_1,J_3)=(2,6)$. Of the remaining nine types, one has $(J_1,J_3)=(6,2)$ and eight have $J_1=J_3\in\{0,1,2\}$. For the growing type,
\begin{itemize}[leftmargin=2em]
\item $\operatorname{supp}w_1=\{J\equiv\pm2\pmod{16}\}$ and $\operatorname{supp}w_3=\{J\equiv\pm6\pmod{16}\}$;
\item all nonzero entries of $w_1$ and $w_3$ have the same absolute value.
\end{itemize}
\end{proposition}

\begin{proof}
Exact computation in $\Z[\zeta_{128}]$. The absolute values were compared through the exact products $w_{a,c}\overline{w_{a,c}}$.
\end{proof}

\begin{lemma}[growing cusps]\label{lem:growing}
Suppose $h/k$ is of growing type and put $t=|q'|=\e^{-2\pi\Rea(1/z)}$. Then
\[
Q_8(\tau)^\delta=\Phi\,\e^{\pi\delta z/k^2}\,\e^{\pi\delta/z}\Big(1-\delta\beta\,\e^{-2\pi/z}+\cR(z)\Big),\qquad|\Phi|=|\beta|=1,
\]
with $|\cR(z)|\le\widetilde F_\delta(t)-1-\delta t$, where
\[
\widetilde F_\delta=(1-\bar U)^{-\delta}(1-\bar V)^{-\delta},\qquad \bar U(t)=\sum_{\substack{J\equiv\pm2\,(16)\\J>2}}t^{(J^2-4)/64},\qquad \bar V(t)=\sum_{\substack{J\equiv\pm6\,(16)\\J>6}}t^{(J^2-36)/64}.
\]
\end{lemma}

\begin{proof}
By Proposition~\ref{prop:types},
\[
R=C\,q'^{-1/2}\,\frac{1+u}{1+v},\qquad u=\sum_{J>2}\alpha_Jq'^{(J^2-4)/64},\qquad v=\sum_{J>6}\beta_Jq'^{(J^2-36)/64},
\]
with $|C|=|\alpha_J|=|\beta_J|=1$. The exponents are integers; $u=O(q'^3)$ and $v=\beta_{10}q'+O(q'^7)$. Put $\beta=\beta_{10}$. From \eqref{eq:theta-quotient}, \eqref{eq:tauprime} and \eqref{eq:R},
\[
e(-\delta\tau/2)=e(-\delta h/2k)\,\e^{\pi\delta z/k^2},\qquad q'^{-\delta/2}=e(\delta d/2k)\,\e^{\pi\delta/z},\qquad q'=e(-d/k)\,\e^{-2\pi/z}.
\]
All unimodular constants are absorbed into $\Phi$ and $\beta$. For the $\delta$-th power: on the connected set $\{\Rea z>0,\ \Rea(1/z)\ge1\}$ we have $|u|,|v|<\frac12$, so $\delta\log Q_8$ and $\delta\big(\text{logarithm of the displayed product}\big)$, built from principal logarithms of $1+u$ and $1+v$, differ by a continuous function with values in $2\pi i\delta\Z$, hence by a constant; this constant is absorbed into $\Phi$. Expand $(1+u)^\delta(1+v)^{-\delta}$ binomially and use $|\binom{\pm\delta}{r}|\le\binom{\delta+r-1}{r}$ for $\delta>0$. Every coefficient is then majorized by the corresponding coefficient of $\widetilde F_\delta(t)$. The coefficient of $t$ in $\widetilde F_\delta$ is $\delta$, coming from $\beta_{10}$, and all other exponents are integers $\ge2$.
\end{proof}

\begin{lemma}[which cusps grow]\label{lem:which}
If $h/k$ is of growing type, then $4\mid k$, so for each $k$ at most $\varphi(k)\le k/2$ fractions $h/k$ are growing. For $k\le40$, $h/k$ is growing if and only if $8\mid k$ and $h\equiv3,5\pmod8$.
\end{lemma}

\begin{proof}
Lemma~\ref{lem:poisson} and \eqref{eq:R} give the same expansion of $S_a$ in powers of $X=|q'|^{1/64}\cdot(\text{phase})$, so by uniqueness of such expansions $J_a=\min\{\ell\ge0:\cG_a(\ell)+\cG_a(-\ell)\ne0\}$. By Proposition~\ref{prop:types}, a growing cusp has $J_1=2$, so $\cG_1(2)\ne0$ or $\cG_1(-2)\ne0$.

Write $s=r+8m$ with $r=a-4$ and $m\bmod2k$. Then $\psi_a(s)=(-1)^m=e(km/2k)$ and
\[
\cG_a(\ell)=e\Big(\frac{hr^2+\ell r}{16k}\Big)\,g(8h,\,2hr+\ell+k;\,2k),\qquad g(A,B;K)=\sum_{m\bmod K}e\Big(\frac{Am^2+Bm}K\Big).
\]
We use two elementary facts.
\begin{enumerate}[label=(\roman*),leftmargin=2em]
\item Let $g_0=\gcd(A,K)$, $A=g_0A'$, $K=g_0K'$. Writing $m=m_0+K'j$ with $m_0\bmod K'$ and $j\bmod g_0$, we have $Am^2/K\equiv A'm_0^2/K'\pmod1$, so
\[
g=\sum_{m_0}e\Big(\frac{A'm_0^2}{K'}+\frac{Bm_0}K\Big)\sum_{j\bmod g_0}e(Bj/g_0).
\]
Hence $g=0$ unless $g_0\mid B$, and in that case $g=g_0\,g(A',B/g_0;K')$.
\item If $\gcd(A,K)=1$, substituting $m=m'+t$ gives
\[
|g|^2=\sum_te\Big(\frac{At^2+Bt}K\Big)\sum_{m'}e\Big(\frac{2Atm'}K\Big)=K\big(1+[2\mid K]\,e(AK/4+B/2)\big).
\]
\end{enumerate}
Take $A=8h$, $K=2k$, $B=2hr+\ell+k$. If $k$ is odd then $g_0=2$, and $2\mid B$ forces $\ell$ to be odd. If $k\equiv2\pmod4$ then $h$ is odd and $g_0=4$; since $2hr\equiv2$ and $k\equiv2\pmod4$, the condition $4\mid B$ forces $\ell\equiv0\pmod4$. In both cases $\cG_1(\pm2)=0$, a contradiction. The bound $\varphi(k)\le k/2$ holds because $k$ is even.

The final statement was verified for all $490$ fractions with $k\le40$ by exact zero tests of $\cG_a(\ell)+\cG_a(-\ell)$ modulo the cyclotomic polynomial $\Phi_{16k}$.
\end{proof}

\begin{proposition}[non-growing cusps]\label{prop:KNG}
For every non-growing type and every $\tau'$ with $\Ima\tau'\ge1$ we have $|R(\tau')|\le\KNG:=23.6475$. Consequently, at every non-growing cusp and for every $z$ with $\Rea(1/z)\ge1$,
\[
|Q_8(\tau)^\delta|\le\e^{\pi\delta\Rea z/k^2}\,\KNG^{\,\delta}\qquad(\delta>0).
\]
\end{proposition}

\begin{proof}
Each $R$ is periodic with period equal to the length $L\in\{2,4,8\}$ of its cycle, so it suffices to consider $0\le\Rea\tau'\le L$.
\begin{itemize}[leftmargin=2em]
\item On $\Ima\tau'\ge6$ we use the bound
\[
|R|\le\frac{|A_{1,J_1}|}{|A_{3,J_3}|}\,t^{(J_1^2-J_3^2)/64}\,\frac{1+\varepsilon_1}{1-\varepsilon_3},\qquad \varepsilon_a=\sum_{J>J_a}\frac{|A_{a,J}|}{|A_{a,J_a}|}\,t^{(J^2-J_a^2)/64},
\]
with $t=\e^{-2\pi\Ima\tau'}$. It is non-increasing in $\Ima\tau'$ because $J_1\ge J_3$, so it suffices to evaluate it at $\Ima\tau'=6$.
\item The rectangle $[0,L]\times[1,6]$ is covered by complex balls on which the theta series, truncated at $J\le40$ with an explicit error term $2\max|A|\,\e^{-2\pi\Ima\tau'\cdot41^2/64}$, are evaluated in ball arithmetic. The rectangles are bisected best-first until every certified upper bound is within a factor $1.02$ of the largest sampled value.
\end{itemize}
The largest certified bound is $23.64743\ldots$, while the largest sampled value is $23.184$. The bound for $Q_8$ follows from \eqref{eq:theta-quotient}, \eqref{eq:tauprime} and $|e(-\delta\tau/2)|=\e^{\pi\delta\Rea z/k^2}$.
\end{proof}

\subsection{Exact data at the cusps of denominators 8 and 16}

\begin{lemma}\label{lem:exact}
For every $n$ the following hold.
\begin{enumerate}[label=(\roman*),leftmargin=2em]
\item The main-term coefficients $\Phi\,e(-nh/k)$ of Lemma~\ref{lem:growing} sum to $2\cos(3\pi n/4)$ over $h/k\in\{3/8,5/8\}$, and to $4\cos(\pi n/8)\cos(\pi(n+\delta)/2)$ over $h/k\in\{3/16,5/16,11/16,13/16\}$.
\item The correction coefficients $\Phi\beta\,e(-nh/k)$ sum to $2\cos\big((3n-1)\pi/4\big)$ over $h/k\in\{3/8,5/8\}$. Consequently the total correction term at $k=8$ in Proposition~\ref{prop:expansion} is $-2\delta\cos\big((3n-1)\pi/4\big)\cP_8(\delta-2)$.
\end{enumerate}
\end{lemma}

\begin{proof}
Put $c_{h,k}=\widehat\cG_1(4)/\widehat\cG_3(36)$, where $\widehat\cG_a(e)=\sum_{\ell^2=e}\cG_a(\ell)$. Exact computation in $\Z[\zeta_{128}]$ and $\Z[\zeta_{256}]$ gives $c_{h,k}=\zeta_{16k}^{t}$ with
\[
t=24,\,40\quad(h/k=3/8,\,5/8),\qquad t=216,\,232,\,152,\,168\quad(h/k=3/16,\,5/16,\,11/16,\,13/16).
\]
It also gives $\widehat\cG_3(100)/\widehat\cG_3(36)=\rho_h$ with $\rho_3=e(1/8)$ and $\rho_5=e(-1/8)$ at $k=8$. The branch integers $m_{h,k}$ of Lemma~\ref{lem:branch} are $0,0,0,0,1,1$ in the order $3/8,5/8,3/16,5/16,11/16,13/16$. Hence the $\delta$-dependent phase of the main term,
\[
\delta\big(\Arg c_{h,k}+2\pi m_{h,k}-\pi h/k\big),
\]
equals $0$ at $k=8$, equals $-\pi\delta/2$ at $h=3,5$ with $k=16$, and equals $+\pi\delta/2$ at $h=11,13$ with $k=16$. Summing complex-conjugate pairs gives (i).

For (ii), by Lemma~\ref{lem:poisson} the first correction in $(S_1/S_3)^\delta$ is $-\delta\rho_hX^{64}=-\delta\rho_h\e^{-2\pi/z}$, so $\beta_{h,8}=\rho_h$, while $\Phi_{h,8}=1$ by (i). Summing over $h=3,5$ gives $e(-3n/8+1/8)+e(-5n/8-1/8)=2\cos\big((3n-1)\pi/4\big)$.
\end{proof}

\begin{lemma}[branch constants]\label{lem:branch}
Let $\mathcal D=\{z:\Rea z>0,\ \Rea(1/z)\ge1\}$ and let $\log Q_8=\sum_{m\ge1}b_mq^m/m$. For each $h/k$ in Lemma~\ref{lem:exact} there is an integer $m_{h,k}$ such that, for all $z\in\mathcal D$,
\[
\log Q_8(\tau)=-\pi i\tau+\Log c_{h,k}-32\log X+\Log(1+u)-\Log(1+v)+2\pi i\,m_{h,k},
\]
where $\log X=-\pi/(32z)$. The integers $m_{h,k}$ are those listed in the proof of Lemma~\ref{lem:exact}.
\end{lemma}

\begin{proof}
The set $\mathcal D$ is connected. On it $t=|q'|\le\e^{-2\pi}$, so by the proof of Lemma~\ref{lem:growing}, $|u|\le\bar U(\e^{-2\pi})<10^{-8}$ and $|v|\le\bar V(\e^{-2\pi})<2\cdot10^{-3}$; hence every principal logarithm on the right is continuous. The difference of the two sides is therefore a continuous function with values in $2\pi i\Z$, hence constant. At $z=1$ it was evaluated in $160$-bit ball arithmetic. The series was summed over $m\le8000$, with the tail bounded using $|b_m|/m\le1+\ln m$, and each term $q^m$ was computed directly as $\e^{2\pi im\tau}$. In each case the result equals the stated integer to within $3\cdot10^{-11}$.
\end{proof}

\section{The explicit circle method}\label{sec:circle}

For $k\ge1$, $B>0$ and $y=\sqrt{2n+\delta}$ put
\[
\cP_k(B)=\frac{2\pi}k\,\frac{\sqrt B}{y}\,I_1\Big(\frac{2\pi}k\sqrt B\,y\Big),\qquad\text{and}\qquad \cP_k(B)=0\ \text{ for } B\le0 .
\]

\begin{lemma}\label{lem:bessel}
Let $K$ be the circle $|z-\frac12|=\frac12$, oriented negatively. Then for $B>0$,
\[
\frac i{k^2}\int_K\exp\Big(\frac{\pi B}z+\frac{\pi(2n+\delta)z}{k^2}\Big)\,dz=\cP_k(B).
\]
\end{lemma}

\begin{proof}
The map $w=1/z$ sends $K$ to the line $\Rea w=1$. The claim follows from the integral representation $I_1(x)=\frac{x/2}{2\pi i}\int_{1-i\infty}^{1+i\infty}\e^{s+x^2/4s}s^{-2}\,ds$ with $s=\pi Bw$.
\end{proof}

\begin{remark}
As a consistency check of Lemma~\ref{lem:bessel}, of the orientation conventions, and of Lemma~\ref{lem:exact}, the sum of the main terms at $k=8$ and $k=16$ reproduces the exact coefficients, e.g.\ $c_{3.2}(488)=6.6426668\cdot10^{16}$ against $6.6426668\cdot10^{16}$, and $c_{3.2}(490)=4.8976498\cdot10^{9}$ against $4.8976517\cdot10^{9}$; the difference is of the size of the omitted terms.
\end{remark}

\begin{lemma}[Rademacher path]\label{lem:farey}
Let $N\ge1$ and let $\frac{h_1}{k_1}<\frac hk<\frac{h_2}{k_2}$ be consecutive in the Farey sequence of order $N$, taken cyclically on $[0,1)$. Let $\gamma_{h,k}$ be the upper arc of the Ford circle $C(h/k)$, which has centre $\frac hk+\frac i{2k^2}$ and radius $\frac1{2k^2}$, between its points of tangency with $C(h_1/k_1)$ and $C(h_2/k_2)$.
\begin{enumerate}[label=(\alph*),leftmargin=2em]
\item The arcs $\gamma_{h,k}$ form a path from some $\tau_0$ to $\tau_0+1$, and
\[
c_\delta(n)=\sum_{k\le N}\sum_h\frac i{k^2}\,e(-nh/k)\int_{z'}^{z''}Q_8\Big(\frac hk+\frac{iz}{k^2}\Big)^\delta\,\e^{2\pi nz/k^2}\,dz,
\]
where each integral runs along the arc of $K$ through $z=1$ from $z'=\frac{k^2+ikk_1}{k^2+k_1^2}$ to $z''=\frac{k^2-ikk_2}{k^2+k_2^2}$.
\item We have $|z'|,|z''|\le\sqrt2\,k/(N+1)$ and $\Rea z',\Rea z''\le2k^2/(N+1)^2$.
\item On the chord $s_{h,k}=[z',z'']$ we have $\Rea(1/z)\ge1$, $0<\Rea z\le2k^2/(N+1)^2$, and $|s_{h,k}|\le2\sqrt2\,k/(N+1)$.
\item On the two arcs of $K$ from $0$ to $z'$ and from $z''$ to $0$ we have $\Rea(1/z)=1$ and $\Rea z\le2k^2/(N+1)^2$, and each has length at most $\frac\pi2\cdot\sqrt2\,k/(N+1)$.
\end{enumerate}
\end{lemma}

\begin{proof}
(a) Adjacent Ford circles are tangent, so consecutive arcs join, and the integrand is holomorphic in $\HH$ and $1$-periodic. The change of variable $z=-ik^2(\tau-h/k)$ maps $C(h/k)$ onto $K$, sending its top point $\frac hk+\frac i{k^2}$ to $z=1$. The point of tangency with $C(h_1/k_1)$ is
\[
\frac hk-\frac{k_1}{k(k^2+k_1^2)}+\frac i{k^2+k_1^2},
\]
which maps to $z'$; similarly for $z''$. Traversing the upper arc from left to right corresponds to going around $K$ clockwise from $z'$ through $1$ to $z''$. Finally $d\tau=\frac i{k^2}\,dz$ and $e(-n\tau)=e(-nh/k)\,\e^{2\pi nz/k^2}$.

(b) Farey neighbours satisfy $k+k_i\ge N+1$, hence $k^2+k_i^2\ge(N+1)^2/2$. Therefore $|z'|=k(k^2+k_1^2)^{-1/2}\le\sqrt2k/(N+1)$ and $\Rea z'=k^2/(k^2+k_1^2)\le2k^2/(N+1)^2$, and likewise for $z''$.

(c) The chord lies in the closed disc bounded by $K$, on which $\Rea(1/z)\ge1$. Along a segment, $\Rea z$ is at most its value at an endpoint, and $|s_{h,k}|\le|z'|+|z''|$.

(d) On $K$ we have $\Rea(1/z)=1$. An arc of $K$ from $0$ to a point $w$ has length at most $\frac\pi2|w|$, and $\Rea z$ is monotone along it.
\end{proof}

If $N\ge\sqrt{4\pi n+2\pi\delta}$, then on all the paths in (c) and (d),
\begin{equation}\label{eq:expbound}
\big|\e^{(2\pi n+\pi\delta)z/k^2}\big|\le\e^{(4\pi n+2\pi\delta)/N^2}\le\e.
\end{equation}

\begin{proposition}[explicit expansion]\label{prop:expansion}
Let $2<\delta\le4$, $n\ge1$ and $N\ge\sqrt{4\pi n+2\pi\delta}$. Then
\[
c_\delta(n)=\sum_{\substack{h/k\ \mathrm{growing}\\ k\le N}}e\Big(-\frac{nh}k\Big)\,\Phi_{h,k}\Big(\cP_k(\delta)-\delta\beta_{h,k}\,\cP_k(\delta-2)\Big)+E_\delta(n),\qquad|\Phi_{h,k}|=|\beta_{h,k}|=1,
\]
and $|E_\delta(n)|\le\cE(\delta,N)$. Here, with $\Sigma_N=(\tfrac12+\tfrac N8)/(N+1)$,
\[
\cE(\delta,N)=\sqrt2\pi\,\Sigma_N\,\e^{1+\pi\delta}(1+\delta\e^{-2\pi})+2\sqrt2\,\Sigma_N\,\e^{1+\pi\delta}\big(\widetilde F_\delta(\e^{-2\pi})-1-\delta\e^{-2\pi}\big)+2\sqrt2\,\e\,\KNG^{\,\delta}.
\]
\end{proposition}

\begin{proof}
We start from Lemma~\ref{lem:farey}(a).

\emph{Growing cusps, main terms.} Insert Lemma~\ref{lem:growing}. The terms $1$ and $-\delta\beta\e^{-2\pi/z}$ produce the exponentials $\exp\big(\pi\delta/z+\pi(2n+\delta)z/k^2\big)$ and $\exp\big(\pi(\delta-2)/z+\pi(2n+\delta)z/k^2\big)$. Completing each integral from the arc $z'\to1\to z''$ to all of $K$ and applying Lemma~\ref{lem:bessel} gives $\cP_k(\delta)$ and $\cP_k(\delta-2)$. By Lemma~\ref{lem:farey}(d) and \eqref{eq:expbound}, the two added arcs cost at most $\frac1{k^2}\cdot2\cdot\frac\pi2\cdot\frac{\sqrt2k}{N+1}\,\e^{1+\pi\delta}(1+\delta\e^{-2\pi})$ per growing cusp.

\emph{Counting.} By Lemma~\ref{lem:which} there are $2$ growing fractions with $k=8$, $4$ with $k=16$, none with other $k\le23$, and at most $k/2$ for each $k\ge24$ with $4\mid k$. There are at most $N/4$ such $k\le N$, so $\sum_{\text{growing}}1/k\le\frac12+\frac N8$. This gives the first term of $\cE$.

\emph{Growing cusps, remainder.} Move the $\cR$-integral to the chord, which is allowed by Cauchy's theorem since the integrand is holomorphic in $\Rea z>0$. Put $w=\Rea(1/z)\ge1$. Then
\[
|\e^{\pi\delta/z}\cR(z)|\le\e^{\pi\delta w}\sum_{j\ge2}\varphi_j\e^{-2\pi jw},\qquad \varphi_j\ge0 .
\]
Each term is non-increasing in $w$ because $\delta\le4$, so the right side is maximal at $w=1$. The chord length from Lemma~\ref{lem:farey}(c) and the same count give the second term.

\emph{Non-growing cusps.} Move the integral to the chord. There Proposition~\ref{prop:KNG} and \eqref{eq:expbound} bound the integrand by $\e\,\KNG^{\,\delta}$. Summing $\frac1{k^2}\cdot\frac{2\sqrt2k}{N+1}$ over at most $\varphi(k)\le k$ values of $h$ and over $k\le N$ gives the third term.
\end{proof}

\begin{remark}\label{rem:E}
The constant $\cE(\delta,N)$ is increasing in $\delta$, because $\widetilde F_\delta=\exp(\delta W)$ where $W$ has non-negative coefficients, and it is decreasing in $N$. Numerically,
\[
\cE(8/3,81)=4.2\cdot10^4,\qquad \cE(3,81)=1.2\cdot10^5,\qquad \cE(4,81)=2.9\cdot10^6 .
\]
\end{remark}

\section{Proof of Theorem~\ref{thm:A}}\label{sec:A}

Fix $\delta\in[8/3,4]$ and $N=\lceil\sqrt{4\pi n+2\pi\delta}\,\rceil$. Using Lemma~\ref{lem:exact} for $k\in\{8,16\}$ and $|\Phi|=|\beta|=1$ for $k\ge24$, Proposition~\ref{prop:expansion} gives
\[
c_\delta(n)=\mathcal M_8+\mathcal C_8+\mathcal T,
\]
where
\[
\mathcal M_8=2\cos\Big(\frac{3\pi n}4\Big)\cP_8(\delta),\qquad \mathcal C_8=-2\delta\cos\Big(\frac{(3n-1)\pi}4\Big)\cP_8(\delta-2),
\]
and
\begin{equation}\label{eq:T}
|\mathcal T|\le\underbrace{4\Big|\cos\frac{\pi n}8\cos\frac{\pi(n+\delta)}2\Big|\,\cP_{16}(\delta)+4\delta\,\cP_{16}(\delta-2)}_{T_{16}}+\underbrace{\sum_{\substack{24\le k\le N\\4\mid k}}\frac k2\big(\cP_k(\delta)+\delta\,\cP_k(\delta-2)\big)}_{T_{\ge24}}+\cE(\delta,N).
\end{equation}

\begin{proposition}[sign criterion]\label{prop:criterion}
Let $\delta\in[8/3,4]$.
\begin{enumerate}[label=(\alph*),leftmargin=2em]
\item If $n\equiv0,1,3,4,5,7\pmod 8$, then $|\mathcal M_8|\ge\sqrt2\,\cP_8(\delta)$ and $\sgn\mathcal M_8=\sigma(n)$. Hence $\sgn c_\delta(n)=\sigma(n)$ whenever $\sqrt2\,\cP_8(\delta)-2\delta\,\cP_8(\delta-2)>T_{16}+T_{\ge24}+\cE$.
\item If $n\equiv2,6\pmod8$, then $\mathcal M_8=0$, $|\mathcal C_8|=\sqrt2\,\delta\,\cP_8(\delta-2)$ and $\sgn\mathcal C_8=\sigma(n)$. Hence $\sgn c_\delta(n)=\sigma(n)$ whenever $\sqrt2\,\delta\,\cP_8(\delta-2)>T_{16}+T_{\ge24}+\cE$.
\end{enumerate}
\end{proposition}

\begin{proof}
This follows directly from the explicit values of the cosines: $2\cos(3\pi n/4)\in\{\pm2,\pm\sqrt2\}$ with sign $\sigma(n)$ on the residues in (a), and on the residues in (b) we have $-2\delta\cos((3n-1)\pi/4)=+\sqrt2\delta$ for $n\equiv2$ and $-\sqrt2\delta$ for $n\equiv6$.
\end{proof}

\subsection{The critical comparison}
Let $n\equiv2,6\pmod 8$. Then $|\cos(\pi n/8)|=\tfrac{\sqrt2}2$ and $|\cos(\pi(n+\delta)/2)|=|\cos(\pi\delta/2)|$, so
\begin{equation}\label{eq:critical}
\frac{4|\cos\frac{\pi n}8\cos\frac{\pi(n+\delta)}2|\,\cP_{16}(\delta)}{|\mathcal C_8|}=\frac{|\cos(\pi\delta/2)|}{\sqrt{\delta(\delta-2)}}\cdot\frac{I_1\big(\frac\pi8\sqrt\delta\,y\big)}{I_1\big(\frac\pi4\sqrt{\delta-2}\,y\big)}.
\end{equation}
For $\delta\ge8/3$ we have $\frac\pi8\sqrt\delta\le\frac\pi4\sqrt{\delta-2}$, with equality exactly at $\delta=8/3$. Since $I_1$ is increasing, the Bessel quotient is then at most $1$. Moreover
\[
\sup_{[8/3,4]}\frac{|\cos(\pi\delta/2)|}{\sqrt{\delta(\delta-2)}}=0.375,
\]
attained at $\delta=8/3$. On $[8/3,3]$ the function is decreasing, and on $[3,4]$ it is below $0.375$; the latter was certified by bisection into $4000$ subintervals.

\subsection{An envelope for \texorpdfstring{$I_1$}{I1}}

\begin{lemma}\label{lem:envelope}
The function $x\mapsto\sqrt{2\pi x}\,\e^{-x}I_1(x)$ is increasing on $(0,\infty)$ and tends to $1$. Consequently $I_1(x)\le U(x):=\e^x/\sqrt{2\pi x}$ for all $x>0$, and $I_1(x)\ge\kappa(x_0)\,U(x)$ for $x\ge x_0$, where $\kappa(x_0)=\sqrt{2\pi x_0}\,\e^{-x_0}I_1(x_0)$.
\end{lemma}

\begin{proof}
Start from Poisson's integral $I_1(x)=\frac x\pi\int_{-1}^1\sqrt{1-t^2}\,\e^{xt}\,dt$ and substitute $s=x(1-t)$:
\[
\sqrt x\,\e^{-x}I_1(x)=\frac1\pi\int_0^{2x}\sqrt s\,\sqrt{2-s/x}\,\e^{-s}\,ds .
\]
The integrand is non-decreasing in $x$ and the interval of integration grows, so the left side is increasing. By monotone convergence its limit is $\frac{\sqrt2}\pi\Gamma(\frac32)=1/\sqrt{2\pi}$.
\end{proof}

\subsection{The range \texorpdfstring{$701\le n\le20000$}{701 <= n <= 20000}}\label{sec:mid}
Divide $[8/3,3]$ and $[3,4]$ into $16$ equal subintervals each. On a subinterval $[a,b]$, every quantity in Proposition~\ref{prop:criterion} is bounded by monotonicity, using endpoint values:
\begin{itemize}[leftmargin=2em]
\item $\cP_k(B)$ is bounded below by taking $\sqrt B$ minimal, $1/y$ at $y=\sqrt{2n+b}$ and the Bessel argument at $y=\sqrt{2n+a}$, and bounded above symmetrically;
\item $\cE$ is evaluated at $\delta=b$, by Remark~\ref{rem:E}, with any admissible $N$; we take $N=\lceil\sqrt{4\pi n+2\pi b}\,\rceil$;
\item $|\cos(\pi\delta/2)|$ is monotone on each of $[8/3,3]$ and $[3,4]$;
\item $T_{\ge24}$ is bounded by $\#\{24\le k\le N:4\mid k\}\cdot\pi\big(\sqrt\delta\,I_1(\tfrac{2\pi}{24}\sqrt\delta\,y)+\delta\sqrt{\delta-2}\,I_1(\tfrac{2\pi}{24}\sqrt{\delta-2}\,y)\big)/y$.
\end{itemize}
The inequalities of Proposition~\ref{prop:criterion} were verified in ball arithmetic at $100$ bits for every $n\in[701,20000]$ and every subinterval, with no exceptions.

\subsection{The range \texorpdfstring{$n>20000$}{n > 20000}}\label{sec:tail}
Now $y\ge200$. Divide every term of \eqref{eq:T} by the dominant term: $|\mathcal C_8|$ in case (b) and $\sqrt2\,\cP_8(\delta)$ in case (a). Apply Lemma~\ref{lem:envelope} with $x_0=\frac\pi8\sqrt{2/3}\cdot200$, for which $\kappa(x_0)=0.99412\ldots$. Every quotient other than \eqref{eq:critical} is then at most $C\,y^pe^{-gy}$ with $p\le\frac32$ and $g>0$, which is decreasing for $y\ge200$. The constants are evaluated at the worst endpoints of $[8/3,4]$, using that $\frac\pi4\sqrt{\delta-2}-\frac\pi8\sqrt\delta$, $\frac\pi4\sqrt{\delta-2}-\frac\pi{12}\sqrt\delta$ and $\frac\pi8\sqrt\delta$ are increasing in $\delta$ while $\frac\pi4(\sqrt\delta-\sqrt{\delta-2})$ is decreasing. At $y=200$, in ball arithmetic:
\begin{center}
\begin{tabular}{lcccc}
\toprule
case & $k=16$ main & $k=16$ correction & $T_{\ge24}$ & $\cE$\\
\midrule
(b) $n\equiv2,6$ & $\le0.375$ & $1.4\cdot10^{-7}$ & $2.2\cdot10^{-8}$ & $1.3\cdot10^{-46}$\\
(a) other $n$ & \multicolumn{2}{c}{$1.5\cdot10^{-43}$} & $8.2\cdot10^{-64}$ & $5.0\cdot10^{-102}$\\
\bottomrule
\end{tabular}
\end{center}
In case (a) the quotient $2\delta\cP_8(\delta-2)/(\sqrt2\cP_8(\delta))$ is at most $5.2\cdot10^{-40}$. The exponential gaps are at least $0.0859$ and $0.118$ in case (b) and at least $\frac\pi4(2-\sqrt2)$ in case (a). In both cases the total is less than $1$ for all $y\ge200$.

\subsection{The range \texorpdfstring{$n\le700$}{n <= 700}}\label{sec:finite}
Each $c_\delta(n)\in\Q[\delta]$ was computed exactly from \eqref{eq:rec}. For each $n\le700$:
\begin{itemize}[leftmargin=2em]
\item the signs at $\delta=8/3$ and at $\delta=4$ were evaluated exactly;
\item the absence of roots in $(8/3,4)$ was certified by Descartes' rule of signs applied to $(1+s)^{\deg}\,c_{8/3+\frac{4/3}{1+s}}(n)$, bisecting whenever the number of sign variations exceeded $1$.
\end{itemize}
For every $n\le700$ and every $\delta\in[8/3,4]$ the sign is $\sigma(n\bmod8)$, the only zero being $c_4(3)=0$. The same procedure applied to $[\beta^+,8/3]$ with $\beta^+=2.664479110226973>\beta$ shows that the signs are $\sigma(n\bmod8)$ for all $n\le700$ and $\delta\in[\beta^+,8/3]$; applied to $[2.6,8/3]$ it finds exactly one root, namely that of $c_\delta(14)$ at $\delta=\beta$.

\begin{proof}[Proof of Theorem~\ref{thm:A}]
Combine \S\ref{sec:finite} for $n\le700$, \S\ref{sec:mid} for $701\le n\le20000$, and \S\ref{sec:tail} for $n>20000$.
\end{proof}

\section{Proof of Theorem~\ref{thm:B}}\label{sec:B}

Let $2<\delta<8/3$ and $n\equiv10\pmod{16}$. Then $\mathcal M_8=0$ and $\mathcal C_8=+\sqrt2\,\delta\,\cP_8(\delta-2)>0$. By Lemma~\ref{lem:exact}, the main term at $k=16$ equals
\[
4\cos\frac{5\pi}4\,\cos\frac{\pi(10+\delta)}2\,\cP_{16}(\delta)=2\sqrt2\cos\Big(\frac{\pi\delta}2\Big)\cP_{16}(\delta),
\]
which is negative. Since $\delta/4>\delta-2$, the quotient \eqref{eq:critical} tends to $+\infty$ as $n\to\infty$. Every other term of Proposition~\ref{prop:expansion} has Bessel argument at most $\max\big(\frac\pi4\sqrt{\delta-2},\frac\pi{12}\sqrt\delta\big)y<\frac\pi8\sqrt\delta\,y$, or is $O(1)$. Hence $c_\delta(n)<0$ for large $n$. The case $n\equiv14\pmod{16}$ is identical with all signs reversed.

For the explicit example, take $\delta=2.66448\in[\beta,8/3)$ and $n=1\,000\,010$. Proposition~\ref{prop:expansion}, with the exact terms at $k=8$ and $k=16$ and absolute bounds for all others, evaluated in $200$-bit ball arithmetic, gives
\[
c_{2.66448}(1\,000\,010)\le-5.739\cdot10^{387}.
\]
This is $13.3\%$ of the $k=16$ main term. Since $\sigma(1\,000\,010\bmod8)=\sigma(2)=+$, this contradicts \eqref{eq:sigma}.

\emph{Independent confirmation.} The direct computation of \S\ref{sec:computations} (\texttt{q8\_lite.py}, $2600$-bit ball arithmetic, no use of \S\S\ref{sec:cusps}--\ref{sec:circle}) evaluates all $c_{2.66448}(n)$, $n\le860\,000$, with every sign certified. It finds $12\,499$ violations of \eqref{eq:sigma}, all in the classes $n\equiv10,14\pmod{16}$; the first is $c_{2.66448}(760\,014)>0$ (with $760\,014\equiv14\pmod{16}$, where $\sigma$ requires $-$), and from there on every index in these two classes violates the pattern. The location agrees with the crossover predicted by the asymptotic expansion.
\qed

\section{The remaining parts: \texorpdfstring{$\delta=\pm2$}{delta = +-2} and \texorpdfstring{$-1<\delta\le\frac{7-\sqrt{73}}2$}{-1 < delta <= (7-sqrt73)/2}}\label{sec:rest}

\subsection{Negative exponents}
For $\delta=-s<0$ we have $Q_8^{-s}=e(s\tau/2)(S_3/S_1)^s$. The roles of the cusp types are exchanged: the \emph{growing} type is now $(J_1,J_3)=(6,2)$ of Proposition~\ref{prop:types}, for which, by exact computation, $\operatorname{supp}w_3=\{J\equiv\pm2\ (16)\}$, $\operatorname{supp}w_1=\{J\equiv\pm6\ (16)\}$, and all nonzero entries have equal modulus. Hence Lemma~\ref{lem:growing} holds verbatim for $R^{-1}$ with $s$ in place of $\delta$. Its correction term $e^{\pi(s-2)/z}$ is bounded for $s\le2$ and is left in the remainder. The argument of Lemma~\ref{lem:which} with $a=3$ shows that growing cusps have $4\mid k$, and an exact check for $k\le40$ gives exactly $8\mid k$, $h\equiv1,7\pmod8$. The certified bound of Proposition~\ref{prop:KNG} for the other types is $\sup|R^{-1}|\le K'_{\rm NG}=23.6475$. The exact data at $k=8,16$ (Gauss sums in $\Z[\zeta_{128}],\Z[\zeta_{256}]$ and certified branch integers) give, with $y=\sqrt{2n-s}$, the amplitudes
\[
2\cos\frac{\pi(n-s)}4\quad(k=8),\qquad 2\cos\frac{\pi(n-s)}8+2\cos\frac{\pi(7n+s)}8\quad(k=16).
\]
The analogue of Proposition~\ref{prop:expansion} holds for $0<s\le2$ with a single main term per growing cusp and error
\[
\cE'(s,N)=\sqrt2\pi\Sigma_N\e^{1+\pi s}+2\sqrt2\,\Sigma_N\e^{1+\pi s}\big(\widetilde F_s(\e^{-2\pi})-1\big)+2\sqrt2\,\e\,K_{\rm NG}'^{\,s},\qquad N\ge\sqrt{4\pi n}+1,
\]
since $\e^{\pi sw}(\widetilde F_s(\e^{-2\pi w})-1)$ is non-increasing in $w$ for $s\le2$. For part (a), $\delta=2$, the same modification (no correction term extracted) is used in Proposition~\ref{prop:expansion}.

\subsection{The degeneration at \texorpdfstring{$\delta=-1$}{delta = -1}}
By a theorem of Andrews and Bressoud \cite{AB} (see also Richmond--Szekeres \cite[Thm.~5.1]{RS}), the coefficients of $Q_8(q)^{-1}=(q^3,q^5;q^8)_\infty/(q,q^7;q^8)_\infty$ vanish at all $n\equiv3\pmod4$. Consequently, at these $n$ the leading amplitude $2\cos(\pi(n-s)/4)=\mp2\sin(\pi\varepsilon/4)$, $\varepsilon=1-s=\delta+1$, tends to $0$, while the error term does not, and the method of \S\ref{sec:A} alone gives a threshold $n_0(\delta)\to\infty$ as $\delta\to-1$ (this is the obstruction met in \cite{HL}). We remove it by subtracting the expansion at $s=1$:
\[
c_\delta(n)=c_\delta(n)-c_{-1}(n)=\mp2\sin\frac{\pi\varepsilon}4\,\cP_8(s)+\sum_{k\ge16}\big(M_k(s)-M_k(1)\big)+\big(E'(s)-E'(1)\big),
\]
and bounding every difference by $\varepsilon$ times an explicit quantity, using three lemmas.

\begin{lemma}\label{lem:imlog}
For $|q|<1$, $|\Ima\log Q_8(q)|\le\frac{\pi}{2}\frac{1}{1-|q|}$. On every Farey chord $|\Ima\log Q_8|\le0.556\,N^2$, and the phase $\theta_{h,k}$ of the main term at a growing cusp (so that $\Phi_{h,k}(s)=\e^{is\theta_{h,k}}$) satisfies $|\theta_{h,k}|\le k^2/2+0.01$.
\end{lemma}
\begin{proof}
$\log Q_8=\sum_m\epsilon_m\log(1-q^m)$ with principal logarithms and $|\Arg(1-x)|\le\frac\pi2|x|$ for $|x|<1$; sum over odd $m$. On a chord $\Rea z\ge k^2/(2N^2)$, so $1-|q|\ge0.9\pi/N^2$. For $\theta_{h,k}$ evaluate the branch identity of Lemma~\ref{lem:branch} (for $R^{-1}$) at $z=1$, where $1-|q|\ge0.9\cdot2\pi/k^2$.
\end{proof}

\begin{lemma}[arc lemma]\label{lem:arc}
Let $A_s=\int f_s\,dz$ over one of the arcs of $K$ from $0$ to $z'$ or $z''$, with $f_s=\Phi(s)\exp(\pi s/z+(2\pi n-\pi s)z/k^2)$, $s\in[s_0,1]$, $s_0>0$. Then
\[
|\partial_sA_s|\le\e^{1+\pi s}\Big[\Big(\frac{k^2}2+0.01\Big)\frac{\pi}{\sqrt2}\frac{k}{N+1}+\frac{2+\pi/2}{s_0}+\frac{\pi^2/\sqrt2}{k(N+1)}\Big].
\]
\end{lemma}
\begin{proof}
$\partial_sf_s=(i\theta+\pi w-\pi/(k^2w))f_s$ with $w=1/z=1+iv$, $|v|\ge k_1/k$. The terms $i\theta f_s$ and $\pi f_s/(k^2w)$ are bounded using $|f_s|\le\e^{1+\pi s}$ and the arc length. For $\int\pi w f_s\,dz=-\pi\Phi\,i\e^{\pi s}\int\e^{i\pi sv}g(v)\,dv$, $g=\e^{b/(1+iv)}/(1+iv)$, integrate by parts: $|g|\le\e$, and $\int|g'|\le\e\int\big(b(1+v^2)^{-3/2}+(1+v^2)^{-1}\big)dv\le\e(1+\pi/2)$ because $b/(1+v_1^2)=b\Rea z'\le1$.
\end{proof}

\begin{lemma}\label{lem:dP}
For $0<s_0\le s\le1$ and $y=\sqrt{2n-s}$, $\cP_k(s)$ is increasing in $s$ and $\cP_k(1)-\cP_k(s)\le(1-s)D_k\cP_k(1)$, where $D_k=\frac1{2s_0}+\frac1{2y^2}+\frac{a_ky}{2\sqrt{s_0}}\big(1+\frac1{a_k\sqrt{s_0}y}\big)$, $a_k=2\pi/k$.
\end{lemma}
\begin{proof}
$I_0=I_2+2I_1/x\le I_1+2I_1/x$ gives $0<I_1'/I_1=I_0/I_1-1/x\le1+1/x$; differentiate $\log\cP_k$.
\end{proof}

\begin{lemma}[derivatives of the remaining error terms]\label{lem:derr}
Let $s\in[s_0,1]$ and $N\ge\sqrt{4\pi n}+1$.
\begin{enumerate}[label=(\roman*),leftmargin=2em]
\item \emph{Growing cusps, chord remainder.} Write $Q_8^{-s}=\Phi(s)\e^{\pi s/z-\pi sz/k^2}W^s$ with $W=(1+u)/(1+v)$ as in Lemma~\ref{lem:growing} (for $R^{-1}$). Then on the chord $s_{h,k}$,
\[
\big|\partial_s\big[\Phi(s)\e^{\pi s/z-\pi sz/k^2}(W^s-1)\big]\big|\le\e^{\pi s}\Big[\Big(\frac{k^2}2+\frac{2\pi N^2}{k^2}+0.1\Big)\big(\widetilde F_1(\e^{-2\pi})-1\big)+\widetilde F_1(\e^{-2\pi})\,\ell(\e^{-2\pi})\Big],
\]
where $\ell=-\log(1-\bar U)-\log(1-\bar V)$.
\item \emph{Non-growing cusps.} On the chord, $|\partial_sQ_8^{-s}|\le K'_{\rm NG}\big(\log K'_{\rm NG}+0.01+0.556N^2\big)$.
\end{enumerate}
\end{lemma}
\begin{proof}
(i) The derivative equals $(i\theta+\pi/z-\pi z/k^2)\Phi\e^{\cdots}(W^s-1)+\Phi\e^{\cdots}W^s\log W$. On the chord, $|\theta|\le k^2/2+0.01$ by Lemma~\ref{lem:imlog}; $|1/z|\le1/\Rea z\le2N^2/k^2$ since $\Rea z\ge k^2/(2N^2)$; and $\pi|z|/k^2\le0.1$. By the coefficientwise majorization of Lemma~\ref{lem:growing}, $|W^s-1|\le\widetilde F_s(t)-1\le\widetilde F_1(t)-1$, $|W|^s\le\widetilde F_1(t)$ and $|\log W|\le\ell(t)$, with $t=\e^{-2\pi w}$, $w=\Rea(1/z)\ge1$. Both $\e^{\pi sw}(\widetilde F_1(t)-1)$ and $\e^{\pi sw}\widetilde F_1(t)\ell(t)$ are sums of terms $\e^{(\pi s-2\pi j)w}$ with $j\ge1$ and non-negative coefficients, hence non-increasing in $w$ since $s\le1$, so their maxima are at $w=1$.

(ii) $\partial_sQ_8^{-s}=-\log Q_8\cdot Q_8^{-s}$. On the chord $|Q_8^{-s}|\le K'^{\,s}_{\rm NG}\le K'_{\rm NG}$, while $|\Rea\log Q_8|=|\log|Q_8||\le\log K'_{\rm NG}+\pi\Rea z/k^2$, using both certified bounds $|R|\le K_{\rm NG}$ and $|R^{-1}|\le K'_{\rm NG}$ (which are equal to $23.6475$). Finally $|\Ima\log Q_8|\le0.556N^2$ by Lemma~\ref{lem:imlog}.
\end{proof}

Summing Lemma~\ref{lem:derr} over the chords, Lemma~\ref{lem:arc} over the arcs (with the counts of Proposition~\ref{prop:expansion}), Lemma~\ref{lem:dP} and $|\Phi_{h,k}(s)-\Phi_{h,k}(1)|\le\varepsilon|\theta_{h,k}|$ for the main terms with $k\ge24$, and $|A_{16}(n,s)|\le\frac\pi2\varepsilon$ for the $k=16$ amplitude, which vanishes at $s=1$ and has $s$-derivative at most $\pi/2$, we obtain $|c_\delta(n)\mp2\sin(\pi\varepsilon/4)\cP_8(s)|\le\varepsilon\cdot B(n)$ with $B(n)$ explicit and independent of $\varepsilon$.

Dividing by $\varepsilon$ yields a sufficient condition for the sign at $n\equiv3,7\pmod8$ that \emph{does not involve $\varepsilon$}:
\[
\frac\pi2\Big(1-\frac{(\pi\varepsilon_{\max}/4)^2}6\Big)\cP_8(s_0)>\frac\pi2\cP_{16}(1)+\sum_{\substack{24\le k\le N\\4\mid k}}\frac k2\Big(\frac{k^2}2+0.01+D_k\Big)\cP_k(1)+(\text{arc})+(\text{chord})+(\text{non-growing}).
\]
At the other residues $|2\cos(\pi(n-s)/4)|\ge2\cdot0.5698$ for $s\in[0.772,1]$ and the standard criterion applies.

\subsection{Certified verification}
\begin{center}
\begin{tabular}{llll}
\toprule
part & $n\le n_1$ (exact) & $n_1<n\le20000$ (ball arithmetic) & $n>20000$\\
\midrule
(a) $\delta=2$ & integers, $n\le400$ & criterion holds for $n\ge295$ & ratio $\le8.5\cdot10^{-13}$\\
(d) $\delta=-2$ & integers, $n\le400$ & criterion holds for $n\ge297$ & ratio $\le8.5\cdot10^{-13}$\\
(c) $s\in[0.772,1]$ & $\Q[\delta]$ on $[-1,-0.772]$, $n\le400$ & criterion holds for $n\ge372$ & ratio $\le4.3\cdot10^{-26}$\\
\bottomrule
\end{tabular}
\end{center}
For (a) and (d), the dominant term at each residue class modulo $16$ is the $k=8$ term, or the $k=16$ term where the $k=8$ amplitude vanishes ($n\equiv2,6$ for $\delta=2$; $n\equiv0,4\pmod8$ for $\delta=-2$), with the signs of the stated patterns; the exact coefficients for $n\le400$ follow the patterns, with the single zero $c_{-2}(8)=0$. For (c), the finite check divides $c_\delta(n)$ by $\delta+1$ at the $100$ indices $n\equiv3\pmod4$ where $c_{-1}(n)=0$, and $c_\delta(4)$ additionally by $\delta^2-7\delta-6=(\delta-b)(\delta-b')$, which is non-negative on $[-1,b]$; Descartes' rule with bisection then certifies the signs on the closed interval. This proves Theorem~\ref{thm:C}. \qed

\section{Computations}\label{sec:computations}

\subsection*{What the computations rely on}
Every computational step belongs to one of three categories.
\begin{enumerate}[leftmargin=2em]
\item \emph{Exact integer and rational arithmetic}, in Python integers and FLINT \cite{flint} polynomials: Propositions~\ref{prop:orbit} and~\ref{prop:types}, Lemma~\ref{lem:which}, the Gauss sums in Lemma~\ref{lem:exact}, and \S\ref{sec:finite}.
\item \emph{Ball arithmetic}, in Arb \cite{arb} as provided by python-flint~0.9.0 with FLINT~3.6.0 and Python~3.12.3: Proposition~\ref{prop:KNG}, Lemma~\ref{lem:branch}, \S\S\ref{sec:mid}--\ref{sec:tail}, and \S\ref{sec:B}. Arb guarantees that each computed ball contains the exact value, and an inequality is accepted only when the whole ball satisfies it.
\item \emph{Floating point}, used only for exploration and as hash keys in the orbit computation. Every match found through a hash key was confirmed exactly, and no conclusion depends on floating-point arithmetic.
\end{enumerate}

\subsection*{Scripts}
\begin{center}
\begin{tabular}{ll}
\toprule
statement & script\\
\midrule
Lemma~\ref{lem:weil} (numerical sanity check) & \texttt{weil\_check.py}\\
Propositions~\ref{prop:orbit}, \ref{prop:types} & \texttt{orbit.py}, \texttt{cusp\_types.py}\\
Lemma~\ref{lem:which} & \texttt{cusp\_exact.py}\\
Proposition~\ref{prop:KNG} & \texttt{supK.py}\\
Lemma~\ref{lem:exact} & \texttt{exact.py}\\
Lemma~\ref{lem:branch} & \texttt{branch\_cert.py}\\
Remark~\ref{rem:E} & \texttt{consts.py}\\
\S\ref{sec:mid} & \texttt{certify\_mid.py}\\
\S\ref{sec:tail} & \texttt{tail\_sc.py}\\
\S\ref{sec:finite} & \texttt{c3.py}\\
\S\ref{sec:B} & \texttt{discert\_sc.py}\\
Remark ($[\beta^+,8/3]$, $n\le700$) & \texttt{c3\_betaplus.py}\\
\S\ref{sec:rest}: negative-exponent cusp data & \texttt{inv\_exact.py}, \texttt{cusp\_exact\_inv.py}, \texttt{supKinv.py}\\
\S\ref{sec:rest}: part (c) & \texttt{neg\_certify.py}, \texttt{neg\_tail772.py}, \texttt{neg\_finite2.py}\\
\S\ref{sec:rest}: parts (a), (d) & \texttt{pm2\_certify.py}, \texttt{pm2\_tail.py}, \texttt{pm2\_finite.py}\\
independent direct check of Theorem~\ref{thm:B} & \texttt{q8\_lite.py} (or \texttt{q8\_signcheck.py})\\
\bottomrule
\end{tabular}
\end{center}
The script \texttt{q8\_lite.py} computes the coefficients of $Q_8^\delta$ directly from \eqref{eq:rec}, by a divide-and-conquer evaluation with blocked ball-arithmetic convolutions, and does not use \S\S\ref{sec:cusps}--\ref{sec:circle}. Its ball radii grow like $\e^{2.5\sqrt n}$, so $2600$ bits suffice at $n\approx8.6\times10^5$, with a memory footprint below $1$ GB. The older \texttt{q8\_signcheck.py} uses the Newton exponential, whose radii grow like $\e^{6.4\sqrt n}$, and needs $15$--$20$ GB.

\begin{thebibliography}{9}
\bibitem{flint} The FLINT team, \emph{FLINT: Fast Library for Number Theory}, version 3.6.0, \url{https://flintlib.org}.
\bibitem{HL} B. He and L. Li, \emph{Some conjectures of Schlosser and Zhou on sign patterns of the coefficients of infinite products}, arXiv:2509.10023.
\bibitem{arb} F. Johansson, \emph{Arb: efficient arbitrary-precision midpoint-radius interval arithmetic}, IEEE Trans. Comput. \textbf{66} (2017), 1281--1292.
\bibitem{Rad} H. Rademacher, \emph{Topics in Analytic Number Theory}, Grundlehren Math. Wiss. 169, Springer, 1973.
\bibitem{AB} G. E. Andrews and D. M. Bressoud, \emph{Vanishing coefficients in infinite product expansions}, J. Austral. Math. Soc. Ser. A \textbf{27} (1979), 199--202.
\bibitem{BHK} K. Bringmann, B. Heim and B. Kane, \emph{On a sign-change conjecture of Schlosser and Zhou}, J. Math. Anal. Appl. \textbf{548} (2025).
\bibitem{Chern} S. Chern, \emph{Asymptotics for the Fourier coefficients of eta-quotients}, J. Number Theory \textbf{199} (2019), 168--191.
\bibitem{Hirschhorn} M. Hirschhorn, \emph{On the expansion of a continued fraction of Gordon}, Ramanujan J. \textbf{5} (2001), 369--375.
\bibitem{Iseki} S. Iseki, \emph{A partition function with some congruence condition}, Amer. J. Math. \textbf{81} (1959), 939--961.
\bibitem{RS} B. Richmond and G. Szekeres, \emph{The Taylor coefficients of certain infinite products}, Acta Sci. Math. (Szeged) \textbf{40} (1978), 347--369.
\bibitem{SZ} M. J. Schlosser and N. H. Zhou, \emph{On the infinite Borwein product raised to a positive real power}, Ramanujan J. \textbf{61} (2023), 515--543, doi:10.1007/s11139-021-00519-3; arXiv:2011.10552.
\end{thebibliography}

\end{document}
